\color{unidarkgreen}{\subsection{Дифференциально-фазная защита (ДФЗ)}}
\color{black}

\begin{enumerate}

%%%%%%%%%%%%%%%%%%%%%%%%%%%%%%%%%%%%%%%%%%%%%%%%%% ОБШИЕ СВЕДЕНИЯ %%%%%%%%%%%%%%%%%%%%%%%%%%%%%%%%%%%%%%%%%%%%%%%%%%
\item Общие сведения

\begin{enumerate}

\item
Дифференциально-фазная защита (ДФЗ) линии является быстродействующей защитой с абсолютной селективностью и выполняет функцию основной защиты линии от всех видов КЗ.

\item
ДФЗ состоит из полукомплектов, устанавливаемых на концах защищаемой линии. 
Каждый полукомплект включает микропроцессорный терминал защиты и соответствующее высокочастотное (ВЧ) оборудование, обеспечивающее обмен ВЧ-сигналами между полукомплектами по ВЧ-каналу связи с использованием внешних приёмопередатчиков (ПП).

\item 
ДФЗ селективно срабатывает:
\begin{itemize}
    \item при всех видах КЗ, в том числе при постановке линии под напряжение; 
    \item при переходе внешнего КЗ во внутреннее, в том числе при реверсе мощности.
\end{itemize}

\item 
Обеспечивается несрабатывание защиты:
\begin{itemize} 
\item при постановке линии под напряжение и включении в транзит при отсутствии КЗ;
\item при неполнофазных режимах;
\item при внешних КЗ, в том числе при реверсе мощности;
\item при синхронных качаниях и в режиме асинхронного хода;
\item при КЗ за трансформаторами ответвлений (отпаек) на линиях с ответвлениями;
\item при бросках тока намагничивания трансформаторов без КЗ;
\item при несимметрии токов, вызванной тяговой нагрузкой;
\item при внешних КЗ с насыщением трансформаторов тока (при соблюдении заявленных производителем требований к ТТ);
\item при каскадных отключениях на обходных связях, несинхронных включениях и режиме одностороннего включения без КЗ.
\end{itemize}

\item
Принцип действия защиты основан на сравнении фаз токов по обоим концам защищаемой линии. 
Передача фазовой информации с одного конца защищаемой линии на другой осуществляется посредством токов высокой частоты, генерируемых ВЧ-передатчиком. 
Управляющее воздействие на ВЧ-передатчик формируется органом манипуляции на основе тока манипуляции, получаемого с выхода комбинированного фильтра токов. 
Ток манипуляции рассчитывается как линейная комбинация токов прямой и обратной последовательностей по формуле:
\begin{align}
i_{\textmd{ман}} = i_1 + k_{\textmd{ман}} \cdot i_2,
\end{align}
\begin{eqexpl}
    \item{$ i_{\textmd{ман}} $} мгновенное значение тока манипуляции, А;
	\item{$ i_1 $} мгновенное значение тока прямой последовательности, А;
	\item{$ i_2 $} мгновенное значение тока обратной последовательности, А;
	\item{$ k_{\textmd{ман}} $} коэффициент манипуляции, учитывающий отношение тока обратной последовательности к прямой последовательности. Задается уставкой \textbf{«Kман.»}.
\end{eqexpl}

\medskip

\item 
ВЧ-передатчик, управляемый органом манипуляции, генерирует токи высокой частоты пакетами, длительность которых приблизительно равна интервалам перехода тока манипуляции через ноль, т.е. с периодичностью, равной половине периода промышленной частоты. 
Таким образом, фаза этих ВЧ-пакетов соответствует фазе тока манипуляции. 
Орган манипуляции управляет специальным быстродействующим реле на основании выходного сигнала \textit{«ВЧ ПРД»}, которое и отвечает за пуск ВЧ-передатчика.

\item
По умолчанию пуск ВЧ-сигнала осуществляется при отрицательной полуволне тока манипуляции. 
При необходимости можно предусмотреть пуск ВЧ-сигнала при положительной полуволне с помощью программного ключа \textbf{«Инверс. ОМ»}.

\item
Для изменения ширины импульса выходного сигнала органа манипуляции реализован сдвиг тока манипуляции относительно оси абсцисс. 
Сдвиг тока манипуляции задается уставкой \textbf{«Сдвиг ман.»} в диапазоне ±10\% от $ I_{\textmd{ном}} $ относительно нуля.

\item
Для синхронизации фронтов ВЧ-сигналов и обеспечения совместимости различных полукомплектов ДФЗ предусматривается доворот тока манипуляции на заданный угол уставкой \textbf{«Доворот ОМ»}.

\item
\textls[-15]{Функциональная схема модуля управления ВЧ приёмопередатчиком представлена на рисунке \ref{pict_dfz:Modul_manag_VChPP}.}

\item 
\textls[-15]{Уставочные параметры модуля управления ВЧ приёмопередатчиком представлены в таблице \ref{app_set:vchpp}.}

\medskip

\begin{figure}[H]
\centering
\includegraphics[width=1\textwidth,height=1\textheight,keepaspectratio]{\fbpath/01.01 ВЧ/561.0761 ДФЗ/img/Modul_manag_VChPP.pdf}
\captionof{figure}{Модуль управления ВЧПП}\label{pict_dfz:Modul_manag_VChPP}
\end{figure}

\medskip

\item
ДФЗ содержит две основные группы измерительных органов (ИО): блокирующие – разрешающие работу органа манипуляции, и отключающие – отвечающие за пуск органа сравнения фаз (ОСФ). 
В начальный момент повреждения срабатывают блокирующие органы, обеспечивая ускоренный запуск ВЧ-передатчика и передачу информации о фазе тока манипуляции на противоположный конец линии. 
Место КЗ, внутри или вне защищаемой линии, определяется органом сравнения фаз по длительности пауз в ВЧ-сигнале, соответствующих углу между токами по концам линии. 

\item
Разрешение на выдачу манипулированного ВЧ-сигнала (логический сигнал ФСУ \textit{«Разрешение ОМ»}) формируется от блокирующих органов, отстроенных от нагрузочного режима. 
Блокирующие органы реализованы с возможностью независимой настройки их уставок по:
\begin{itemize} 
\item линейному току (модулю разности фазных токов) \textit{«ИО Iл блок ДФЗ»};
\item приращению тока прямой последовательности \textit{«ИО dI1 блок ДФЗ»};
\item приращению тока обратной последовательности \textit{«ИО dI2 блок ДФЗ»};
\item току обратной последовательности \textit{«ИО I2 блок ДФЗ»};
\item току нулевой последовательности \textit{«ИО 3I0 блок ДФЗ»}.
\end{itemize}

\item
Подготовка цепи отключения осуществляется от отключающих органов.
Отключающие органы реализованы с возможностью независимой настройки их уставок по:
\begin{itemize} 
\item линейному току (модулю разности фазных токов) \textit{«ИО Iл откл ДФЗ»};
\item приращению тока прямой последовательности \textit{«ИО dI1 откл ДФЗ»};
\item приращению тока обратной последовательности \textit{«ИО dI2 откл ДФЗ»};
\item току обратной последовательности \textit{«ИО I2 откл ДФЗ»};
\item току нулевой последовательности \textit{«ИО 3I0 откл ДФЗ»};
\item по сопротивлению \textit{«РС откл ДФЗ»}.
\end{itemize}

\item
Измерительные органы ДФЗ представлена на рисунке \ref{pict_dfz:IO_DFZ}.

\medskip


\begin{figure}[H]
\centering
\includegraphics[width=0.95\textwidth,height=0.95\textheight,keepaspectratio]{\fbpath/01.01 ВЧ/561.0761 ДФЗ/img/IO_DFZ.pdf}
\captionof{figure}{Измерительные органы ДФЗ}
\label{pict_dfz:IO_DFZ}
\end{figure}

\item 
\textls[-15]{Уставочные параметры измерительных органов ДФЗ представлены в таблице \ref{app_set:dfz}.}

\end{enumerate}


%%%%%%%%%%%%%%%%%%%%%%%%%%%%%%%%%%%%%%%%%%%%%%%%%% БЛОКИРУЮЩИЕ ИО %%%%%%%%%%%%%%%%%%%%%%%%%%%%%%%%%%%%%%%%%%%%%%%%%%
\item Блокирующие измерительные органы

\begin{enumerate}

\item
Срабатывание блокирующих органов подхватывается RS-триггером, который сбрасывается через 0,6 с после их возврата или при появлении сигнала останова ВЧ-передатчика (логический сигнал \textit{«Запрет пуска ВЧ»}).
Это обеспечивает надежное срабатывание защиты противоположного конца линии при каскадном отключении. Сигнал \textit{«Запрет пуска ВЧ»} формируется в следующих случаях:
\begin{itemize} 
\item при срабатывании внутренних защит и УРОВ (входной сигнал ФСУ \textit{«Останов ВЧ»});
\item при приеме внешнего сигнала \textit{«Останов ВЧ внеш.»}.
Останов ВЧ-передатчика от данного сигнала блокируется в случае срабатывания ДЗШ (входной сигнал ФСУ \textit{«Прием Ср. ДЗШ»}), действующего на отключение выключателя защищаемой линии.
Блокировка продлевается на время, равное 200 мс. 
\end{itemize}

\item
Время возврата логического сигнала \textit{«Запрет пуска ВЧ»} составляет 200 мс.

\item
В устройстве также предусматривается возможность ручного пуска ВЧ-передатчика через входной сигнал ФСУ \textit{«Ручной пуск ВЧ»}.
Выдаваемый ВЧ-сигнал при этом является манипулированным.
Ручной пуск разрешается, если в течение последних 600 мс не было срабатываний блокирующих органов или отключающего реле сопротивления.

\item
Кроме того, предусматривается возможность выдачи сплошного ВЧ-сигнала (логический сигнал \textit{«Пуск ВЧ»}) в следующих случаях:
\begin{itemize} 
\item при срабатывании блокировки при сквозных токах через ошиновку, если введен программный ключ \textbf{«БСТО ДФЗ»};
\item от внешней команды пуска ВЧ-передатчика (входной сигнал ФСУ \textit{«Внешний пуск ВЧ»});
\item при обнаружении неисправности канала связи устройством автоматической проверки канала (АПК), если введен программный ключ \textbf{«Блок. ДФЗ от АПК»};
\item при обнаружении неисправности приемопередатчика, если введен программный ключ \textbf{«Блок. ДФЗ от Неиспр. ПП»};
\item при обнаружении неисправности терминала системой самодиагностики, если введен программный ключ \textbf{«Пуск ВЧ при неиспр. терм.»};
\item при выводе терминала или функции ДФЗ из работы, если введен программный ключ \textbf{«Пуск ВЧ при выводе»}.
\end{itemize}

\item
Выдача сплошного ВЧ-сигнала при срабатывании БСТО и от внешней команды пуска дополнительно блокируется при наличии сигнала останова ВЧ-передатчика (логический сигнал \textit{«Запрет пуска ВЧ»}).
А при неисправности терминала, неисправности канала связи или приемопередатчика, а также при выводе терминала или функции ДФЗ из работы пуск ВЧ-передатчика никак не блокируется.

\item
Функциональная схема управления ВЧПП (ДФЗ) представлена на рисунке \ref{pict_dfz:VChPP}.

\medskip

\begin{figure}[H]
\centering
\includegraphics[width=0.95\textwidth,height=0.95\textheight,keepaspectratio]{\fbpath/01.01 ВЧ/561.0761 ДФЗ/img/VChPP.pdf}
\captionof{figure}{Функциональная схема управления ВЧПП (ДФЗ)}
\label{pict_dfz:VChPP}
\end{figure}

\medskip

\item
Устройство АПК, входящее в состав ВЧПП, используется для автоматической проверки ВЧ-канала связи.
К устройству от ВЧПП подводятся сигналы \textit{«Контакт АПК»} и \textit{«Неиспр. ПП»}.
Наличие сигнала \textit{«Контакт АПК»}, свидетельствует об отсутствии неисправности в канале связи.
При выявлении неисправности канала связи формируется логический сигнал \textit{«Сигн. Неиспр. КС»}.
При ручном пуске ВЧ-передатчика, срабатывании блокирующих органов или отключающего реле сопротивления формируется сигнал \textit{«Запрет АПК»}, блокирующий действие АПК.
Действие АПК также может быть выведено оперативно от внешнего сигнала \textit{«Вывод АПК»}.
При обнаружении неисправности приемопередатчика формируется логический сигнал \textit{«Сигн. Неиспр. ПП»}.
При введенных положениях программных ключей \textbf{«Блок. ДФЗ от АПК»} и \textbf{«Блок. ДФЗ от Неиспр. ПП»} дополнительно формируется сигнал \textit{«Блок. от ВЧПП»}, который действует на вывод ДФЗ из работы и пуск ВЧ-передатчика.

\item
Функциональная схема контроля ВЧ-канала связи представлена на рисунке \ref{pict_dfz:Ctrl VCh DFZ}.

\medskip

\begin{figure}[H]
\centering
\includegraphics[width=0.95\textwidth,height=0.95\textheight,keepaspectratio]{\fbpath/01.01 ВЧ/561.0761 ДФЗ/img/Ctrl VCh DFZ.pdf}
\captionof{figure}{Функциональная схема контроля ВЧ-канала связи}\label{pict_dfz:Ctrl VCh DFZ}
\end{figure}

\medskip

\item
Для сигнализации о наличии ВЧ-сигнала в канале связи реализован орган \textit{«ОСФ-Вызов»}, который вводится в работу при отсутствии сигнала \textit{«Запрет АПК»}.
Время срабатывания и возврата \textit{«ОСФ-вызов»} не превышает 20 мс.
Обеспечивается формирование сигнала \textit{«Вызов»}, если сигнал срабатывания \textit{«ОСФ-вызов»} присутствует непрерывно в течение 5 с.

\item 
Уставочные параметры функциональной схемы управления ВЧПП (ДФЗ) и схемы контроля ВЧ-канала связи представлены в таблице \ref{app_set:dfz}.

\end{enumerate}


%%%%%%%%%%%%%%%%%%%%%%%%%%%%%%%%%%%%%%%%%%%%%%%%%% ОТКЛЮЧАЮЩИЕ ИО %%%%%%%%%%%%%%%%%%%%%%%%%%%%%%%%%%%%%%%%%%%%%%%%%%
\item Отключающие измерительные органы

\begin{enumerate}

\item
Отключающий орган по сопротивлению (отключающее РС) необходим для подхвата кратковременного срабатывания отключающих органов и обеспечения корректной работы ДФЗ при симметричных КЗ на защищаемой линии. 
Подхват вводится на время, равное 200 мс. Для сохранения селективности защиты при отключении внешнего КЗ (при неодновременной остановке работы ВЧПП) предусмотрено ограничение длительности срабатывания отключающего РС до 200 мс. 

\item
Характеристика срабатывания отключающего РС представлена на рисунке \ref{pict_dfz:RS_DFZ}.

\medskip

\begin{figure}[H]
\centering
\includegraphics[width=0.4\textwidth,height=0.4\textheight,keepaspectratio]{\fbpath/01.01 ВЧ/561.0761 ДФЗ/img/RS_DFZ.pdf}
\captionof{figure}{ Характеристика срабатывания отключающего РС ДФЗ}\label{pict_dfz:RS_DFZ}
\end{figure}

\medskip

\item
Отключающее РС имеет направленную характеристику срабатывания с охватом начала координат, которая задается уставками \textbf{«Rср РС откл ДФЗ»}, \textbf{«Хср РС откл ДФЗ»} и \textbf{«Ф1 РС откл ДФЗ»}. 
Для отстройки от нагрузки предусмотрен вырез в характеристике срабатывания, который определяется уставками \textbf{«Rнг РС откл ДФЗ»} и \textbf{«Фнг РС откл ДФЗ»}. 
Работа отключающего РС может быть заблокирована при выявлении неисправности в цепях напряжения при введенном программном ключе \textbf{«БНН РС откл ДФЗ»}.

\item
Для предотвращения ложного срабатывания ДФЗ при КЗ за трансформатором ответвления предусмотрен функциональный блок отстройки от повреждений за трансформаторами ответвлений (ОПТО). 
При введенном программном ключе \textbf{«Контр. ДФЗ от ОПТО»} подготовка цепи отключения разрешается только при наличии подтверждения о КЗ на защищаемой линии от ОПТО.

\item
При срабатывании отключающих органов и наличии подтверждения о КЗ на защищаемой линии от ОПТО формируется логический сигнал \textit{«Разрешение ОСФ»}, отвечающий за пуск ОСФ. 
Подключение ОСФ осуществляется с выдержкой времени: 10 мс для линий без ответвлений и 20 мс — при их наличии. 

\item
ОСФ предназначен для определения места повреждения по сдвигу фаз между ВЧ-пакетами, посылаемыми ВЧ-передатчиками обоих концов линии, то есть в конечном счёте по сдвигу фаз токов манипуляции.
Сравнение фаз осуществляется косвенно – по длительности пауз в ВЧ-сигнале (рисунок \ref{pict_dfz:Action_DFZ}).
При внешнем КЗ токи манипуляции находятся в противофазе, ВЧ-передатчики работают в чередующиеся полупериоды промышленной частоты, формируя непрерывный (сплошной) сигнал в ВЧ-канале, при котором ОСФ не срабатывает.
При внутреннем КЗ токи манипуляции на обоих концах совпадают по фазе, ВЧ-передатчики работают синхронно, что приводит к возникновению пауз в ВЧ-сигнале и срабатыванию ОСФ. 

\medskip

\begin{figure}[H]
\centering
\begin{subfigure}{0.42\textwidth}
    \centering
    \includegraphics[width=\textwidth, trim=0 10mm 85mm 0, clip]{\fbpath/01.01 ВЧ/561.0761 ДФЗ/img/Action_DFZ.pdf}
    \caption{а) внутреннее КЗ}
    \label{pict_dfz:Action_DFZ_a}
\end{subfigure}
%\hfill
\begin{subfigure}{0.41\textwidth}
    \centering
    \includegraphics[width=\textwidth, trim=87mm 10mm 0 0, clip]{\fbpath/01.01 ВЧ/561.0761 ДФЗ/img/Action_DFZ.pdf}
    \caption{б) внешнее КЗ}
    \label{pict_dfz:Action_DFZ_b}
\end{subfigure}

\caption{Принцип действия ДФЗ}
\label{pict_dfz:Action_DFZ}
\end{figure}

\medskip

\item
Длительность пауз, в принимаемом ВЧ-сигнале прямо зависит от величины угла сдвига фаз между токами манипуляции: чем меньше угол, тем длиннее паузы.
Уставками ОСФ задается величина сдвига фаз между токами манипуляции, при которой защита блокируется или срабатывает.
Предельное отклонение угла сдвига фаз от 180°, при котором защита ещё блокируется, называется углом блокировки.

\item
Срабатывание ОСФ происходит при превышении разности фаз уставки угла блокировки.
ОСФ выполнен трёхступенчатым, что позволяет реализовать зависимость времени срабатывания защиты от величины сдвига фаз.
Первая ступень ОСФ срабатывает при фиксировании одной паузы в ВЧ-сигнале, превышающей уставку угла блокировки первой ступени.
Вторая ступень – при двух последовательных превышениях уставки второй ступени.
Третья ступень – при трех последовательных превышениях уставки третьей ступени.
Углы блокировки задаются уставками \textbf{«Угол блок. ОСФ-1»}, \textbf{«Угол блок. ОСФ-2»} и \textbf{«Угол блок. ОСФ-3»} соответственно.

\item
Характеристика срабатывания ОСФ представлена на рисунке \ref{pict_dfz:OSF}.

\medskip

\begin{figure}[H]
\centering
\includegraphics[width=0.6\textwidth,height=0.6\textheight,keepaspectratio]{\fbpath/01.01 ВЧ/561.0761 ДФЗ/img/OSF.pdf}
\captionof{figure}{Характеристика срабатывания ОСФ}
\label{pict_dfz:OSF}
\end{figure}

\medskip

\item
Для компенсации расширения зоны блокировки, вызванного особенностями формирования сигнала \textit{«ВЧ ПРМ»} в ВЧПП используется корректировка уставок углов блокировки ОСФ на величину компенсации, задаваемую уставкой \textbf{«Удлинение ВЧ»}.

\item
После срабатывания ОСФ осуществляется действие ДФЗ на отключение выключателя с регулируемой выдержкой времени \textbf{«Тср ДФЗ»}.
При одновременном срабатывании отключающих органов и появлении сигнала останова ВЧ-передатчика (логический сигнал \textit{«Запрет пуска ВЧ»}) ОСФ шунтируется, и защита действует на отключение выключателя без выдержки времени.

\item
При подключении параллельных линий к общим шинам на обоих концах защита может сработать неселективно из-за реверса мощности.
При повреждении на параллельной линии срабатывают блокирующие органы полукомплектов, приводящие к выдаче ВЧ-передатчиком манипулированного ВЧ-сигнала.
Одновременно срабатывают и отключающие органы, запускающие ОСФ.
Однако срабатывания ОСФ и отключения неповрежденной линии не происходит, так как в данном режиме в ВЧ-канале присутствует сплошной сигнал.
При каскадном отключении одного из выключателей параллельной линии возникает реверс мощности.
Во время переходного процесса могут появиться паузы в ВЧ-сигнале, длительность которых достаточна для срабатывания ОСФ, что приведет к отключению неповрежденной линии.
Для исключения такого режима реализована блокировка при реверсе мощности: при одновременном наличии сигнала от отключающих органов и отсутствии срабатывания ОСФ в течение времени, задаваемого уставкой \textbf{«Тср блок. ДФЗ»}, срабатывает блокировка, которая продлевается на время, регулируемое уставкой \textbf{«Тпрод блок. ДФЗ»}.    

\item
Действие ДФЗ на отключение также блокируется при следующих условиях:
\begin{itemize} 
\item при срабатывании блокировки при сквозных токах через ошиновку, которая вводится программным ключом \textbf{«БСТО ДФЗ»};
\item при выводе терминала или функции ДФЗ из работы;
\item при обнаружении неисправности канала связи устройством автоматической проверки канала (АПК), если введен программный ключ \textbf{«Блок. ДФЗ от АПК»};
\item при обнаружении неисправности приемопередатчика, если введен программный ключ \textbf{«Блок. ДФЗ от Неиспр. ПП»};
\item при переводе действия ДФЗ только на сигнал (от внешнего сигнала \textit{«Перевод на сигнал»}).
\end{itemize}

\item
В сетях с тяговой нагрузкой, где наблюдаются значительные небалансы токов обратной последовательности, рекомендуется применять блокирующие и отключающие органы по приращениям токов прямой и обратной последовательностей.

\item
Время срабатывания защиты (без учета срабатывания выходного реле и с учетом задержки на подключение ОСФ) на отключение на двухконцевых линиях при кратности соответствующих величин, равной 3, не более 35 мс.

\item
Функциональная схема ДФЗ представлена на рисунке \ref{pict_dfz:DFZ}.

\medskip

\begin{figure}[H]
\centering
\includegraphics[width=1\textwidth,height=1\textheight,keepaspectratio]{\fbpath/01.01 ВЧ/561.0761 ДФЗ/img/DFZ.pdf}
\captionof{figure}{Функциональная схема ДФЗ}
\label{pict_dfz:DFZ}
\end{figure}

\medskip

\item 
Уставочные параметры функциональной схемы ДФЗ представлены в таблице \ref{app_set:dfz}.

\end{enumerate}

\end{enumerate}